\documentclass[10pt,a4paper]{article}
\usepackage{amsmath,amssymb,amsthm,graphicx}
\usepackage{hyperref}
\usepackage{latexsym}
\usepackage{amssymb}
\usepackage{amsthm}
\usepackage{ulem}
\usepackage{enumitem}
\setlist{itemsep=0.4pt,topsep=0pt}
\parindent 0pt
\parindent 0pt
\parskip 12pt
\newcommand{\HCF}{\mathrm{hcf}}
\newcommand{\ZZ}{\mathbb{Z}}
\newcommand{\rk}{\mathrm{rk}}
\DeclareMathOperator{\len}{length}
\DeclareMathOperator{\bl}{b}
\DeclareMathOperator{\cl}{cl}
\DeclareMathOperator{\di}{d}
\DeclareMathOperator{\Cl}{Cl}
\DeclareMathOperator{\acl}{acl}
\DeclareMathOperator{\tp}{tp}
\DeclareMathOperator{\Th}{Th}
\newcommand{\Cor}{\mathcal{C}_\mathrm{or}}
\newcommand{\Cun}{\mathcal{C}_\mathrm{un}}
\newcommand{\mons}{\mathbb{M}}
\newcommand{\strong}{\mathrm{strong}}
\newcommand{\axioms}{\Delta_{\mathrm{FSTS}}}
\newtheoremstyle{note}% hnamei
{8pt}% hSpace abovei
{14pt}% hSpace belowi
{\normalfont}% hBody fonti
{}% hIndent amounti1
{\bfseries}% hTheorem head fonti
{ }% hPunctuation after theorem headi
{.5em}% hSpace after theorem headi2
{}% hTheorem head spec (can be left empty, meaning ‘normal’)i
\theoremstyle{note}
\newtheorem{thm}{Theorem}[section]
%\newtheorem{theorem}{Theorem}
\newtheorem{lemma}[thm]{Lemma} 
\newtheorem{df}[thm]{Definition}
\newtheorem{rmk}[thm]{Remark} 
\newtheorem{prop}[thm]{Proposition}
\newtheorem{fact}[thm]{Fact}
\newtheorem{cor}[thm]{Corollary}
\newtheorem{conj}[thm]{Conjecture}
\newtheorem{ex}[thm]{Example}
\newtheorem{qu}[thm]{Question}
\newtheorem*{notation}{Notation} 
\title{The theory of free Steiner triple systems is stable}
\author{Silvia Barbina, Enrique Casanovas}
\date{arxiv:2411.13723v3; Creative Commons Attribution 4.0 International}
\begin{document}
\maketitle
\noindent\textit{Editorial note.} The counted statement is the author's environment at source0.tex lines 1033-1034 and its proof environment opens at line 1035; the pairing is by adjacency and the retained proof range is the whole environment.\par\medskip
\noindent\textit{Editorial note.} The preamble is the author's own package lines, macro block and theorem-environment declarations (source lines 4-65); the author's optioned amsmath line and the fleqn-dependent layout line are not retained.\par\medskip
\noindent\textit{Editorial note.} The source carries no label-number dump and no counter values are imposed, so the printed numbering is the wrapper's default; identity is the author label, or the pinned source position where the author gave none.\par\medskip
\noindent\textit{Editorial note.} Metadata provenance: the source record of the pinned copy and the licence evidence file are both present; the licence is CC BY 4.0 with status supported, and the card's source\_version matches the value recorded in that evidence file.\par\medskip
\noindent\textit{Editorial scope.} One original result of the pinned author manuscript, counted once: stability of the theory of free Steiner triple systems (author label stable).\par\medskip
\noindent\textit{Allowed standard prerequisites (editorial).}\begin{itemize}
\item Stability. A theory is stable when over every model the number of types in one variable is at most the cardinality of the model raised to the power omega.. It is the conclusion of the statement.
\item Type over a model. A type is a consistent set of formulas in one free variable with parameters in the model.. It is what the proof counts.
\item Isomorphism type of a countable tuple. Countable tuples over a model are classified up to partial isomorphism by their quantifier-free type in the language of the partial structure.. It is the invariant that bounds the number of types.
\end{itemize}\medskip
	\begin{cor}\label{Feb_23_2024_9}  If $M_1\preccurlyeq M_2$  are unconfined  STSs, then  $M_1\leq M_2$.
	\end{cor}

	\begin{lemma}\label{elementary2}
		Let $M$ be an infinite unconfined STS and let $\bar{a},\bar{b}$ be tuples from $M$. If for some HF-orderings $<$ and $<^\prime$  of $M$,  $\cl_<(\bar{a})$ and $\cl_{<^\prime}(\bar{b})$ have the same diagram as partial STS, then  $\tp_M(\bar{a})=\tp_M(\bar{b})$.
	\end{lemma}

	\begin{prop}\label{stable} The theory of free STSs  is stable.
	\end{prop}

	\begin{proof}  We count  $1$-types over a model $M$ and show that their number is $\leq |M|^\omega$. We assume $M$ is an elementary substructure of the monster model $\mons$. Then  $M\leq \mons$ by  Corollary~\ref{Feb_23_2024_9}, and we can fix an HF-ordering $<$ of $\mons$ where $M$ is an initial segment. If $a,b \in \mons$ and $\cl_<(a)$, $\cl_<(b)$ have the same diagram over $M$ as partial STS, then, by Lemma~\ref{elementary2}, they have same type over $M$.
	Hence, the number of $1$-types over $M$ is bounded by the number of partial STS isomorphism types of countable tuples over $M$, which is clearly $\leq |M|^\omega$.
	\end{proof}

\end{document}
